\section{Experiments}
\label{sec:experiments}

\subsection{Setup, and an honest account of scale}

Our experimental environment cannot reach \texttt{huggingface.co}: the network
policy denies the connection, as it does \texttt{download.pytorch.org}. Public
checkpoints were therefore unavailable and every measurement below is on a model
we trained ourselves. That model satisfies Definition~\ref{def:model} exactly ---
pre-norm residual blocks, RMSNorm, an unembedding --- with $L=6$, $d=128$,
$V=2048$, $1.71$M parameters, trained on a $53$k-token corpus with a held-out
split and early stopping.

We state the consequences plainly rather than in a footnote. The model is
\emph{weak}: validation perplexity is $431$ against a vocabulary of $2048$, i.e.\
about $4.8\times$ better than uniform. Two specific distortions follow. First,
diffuse output distributions overlap more, so measured acceptance rates
$\alpha_\ell = 1 - \TV(p_\ell,p_L)$ are \emph{inflated} relative to a
well-trained model. Second, at $L=6$ the shallowest available exit is
$\rho = 0.167$, whereas a $28$-layer model offers far finer granularity exactly
in the regime where early exit is attractive. Quantitative claims below should
therefore be read as indicative; the qualitative ones
(Sections~\ref{sec:exp-exact} and \ref{sec:exp-cost}) do not depend on scale, and
the two errors we report finding are architecture-independent.

\subsection{Exactness is confirmed at the sampling-noise floor}
\label{sec:exp-exact}

Theorem~\ref{thm:exact} predicts that speculative sampling with \emph{any} draft
returns the target distribution exactly. Testing this against zero would be
meaningless, since any finite sample has positive total variation from $p_L$; the
correct reference is the noise floor, namely $\TV$ between $p_L$ and an equal
number of direct draws from it. Table~\ref{tab:exact} reports both for
$4\times10^{5}$ samples.

\begin{table}[h]
\centering
\begin{tabular}{rrrrrr}
\toprule
layer & $\TV(\text{spec},p_L)$ & $\TV(\text{direct},p_L)$ & ratio
      & $1-\TV(p_\ell,p_L)$ & observed accept \\
\midrule
1 & 0.0202 & 0.0196 & 1.033 & 0.3274 & 0.3271 \\
3 & 0.0202 & 0.0196 & 1.033 & 0.5831 & 0.5833 \\
5 & 0.0186 & 0.0193 & 0.965 & 0.7988 & 0.7978 \\
\bottomrule
\end{tabular}
\caption{Speculative sampling is indistinguishable from direct sampling
(ratio $\approx 1$), and the observed acceptance rate matches the predicted
$1-\TV(p_\ell,p_L)$ to within $10^{-3}$.}
\label{tab:exact}
\end{table}

This is the result that licenses Corollary~\ref{cor:shortlist}, since the
shortlisted draft assigns zero probability to most of the vocabulary and would be
invalid under any support hypothesis.

\subsection{The acceptance bound is directionally right and numerically useless}

Acceptance rises monotonically with depth exactly as tail energy falls:
$\alpha_\ell = 0.327, 0.383, 0.424, 0.520, 0.693, 1.000$ against
$T_\ell = 20.2, 17.8, 15.4, 12.6, 8.1, 0$. The shape predicted by
Theorem~\ref{thm:accept} is therefore present.

The bound itself is not usable. At $\ell=5$, where $\alpha_5 = 0.693$, feeding
the bound the \emph{measured} logit spread gives $e^{-4.21} = 1.5\times10^{-2}$,
already a factor of $46$ too small; the Cauchy--Schwarz step then inflates the
spread by a further $2.5\times$, landing at $\approx 10^{-5}$.
Remark~\ref{rem:loose} anticipated looseness and attributed it primarily to the
$D_U$ step. That attribution was wrong: the softmax stage alone is fatal, and
Cauchy--Schwarz merely compounds it. A useful bound needs a different proof
strategy, not a tighter constant.

\subsection{A cost-model error, caught by a sanity check}
\label{sec:exp-cost}

The most consequential empirical finding is an error in our own accounting. An
earlier version of Proposition~\ref{prop:cost} charged drafting and verification
but omitted the full-depth pass at the bonus position $t+\gamma$. Evaluated at
$\rho = 1$ it returned $S = 1.735$, which is impossible: at $\rho=1$ the draft
\emph{is} the target, so the procedure degenerates to ordinary autoregressive
decoding and $S = 1$ is forced. Restoring the omitted pass gives exactly
$S = 1.000$ (Remark~\ref{rem:sanity}). We recommend the $\rho\to1$ limit as a
standing check on any speculative-decoding cost model; it is cheap and it caught
an error that survived both derivation and adversarial review.

The same experiment overturned our claim about cache reuse. Measured at $w=1$,
reuse is \emph{worse} at every depth ($2.25$ against $1.50$ at $\rho=0.25$,
$\gamma=2$), because splitting verification across two depths pays for the lower
blocks' weight streaming twice; at $w=0$ it helps, $0.929$ against $0.667$. This
is Proposition~\ref{prop:reuse}, and Remark~\ref{rem:kvreuse} retracts the
earlier claim.

\subsection{The head cost decides whether self-drafting pays}

With the corrected model, the best attainable memory-bound speedup on our model
is $S = 1.026$, at $\ell=1$, $\gamma=1$, with head fraction $u = 0.153$. Setting
$u = 0.26$, the value for Qwen3-0.6B, \emph{no} early-exit layer exceeds $S=1$:
the optimum retreats to the degenerate $\rho=1$. The margin is entirely
determined by the unembedding cost, exactly as Corollary~\ref{cor:ceiling}
predicts, and a $2.6\%$ gain on a model whose acceptance rates are already
inflated is not a result we would build on.

One further asymmetry deserves recording. At $\ell=1$ the sampling-mode
acceptance rate is $\alpha_1 = 0.327$ while top-$1$ agreement between $p_1$ and
$p_L$ is only $0.052$. These govern different deployment modes --- the former
speculative sampling at temperature $1$, the latter greedy decoding --- and the
gap of more than $6\times$ means an acceptance rate that looks tolerable for
sampled decoding is close to useless for greedy decoding. Reported acceptance
rates should always name the mode.

\subsection{Yield: acceptance decays faster than geometric}

Running $200$ real speculative rounds at $\ell=1$, $\gamma=4$ gives measured
$\E[Y] = 1.320$ against $1.481$ from the geometric formula of
Corollary~\ref{cor:geom} evaluated at $\bar\alpha = 0.327$. The measured prefix
probabilities are $\Prob[A_j] = 0.250, 0.055, 0.010, 0.005$ against
$\bar\alpha^{j} = 0.327, 0.107, 0.035, 0.011$ under independence: acceptance
decays \emph{faster} than geometric within a round.

This contradicts the remark following Corollary~\ref{cor:geom}, which argued from
Jensen's inequality that cross-context heterogeneity makes the geometric formula
a conservative under-estimate. That argument concerns variation \emph{between}
contexts; within a round the dominant effect runs the other way, because each
accepted draft moves the model onto its own continuation, where the shallow-layer
draft agrees less well. Practitioners estimating speedup from a mean acceptance
rate will therefore \emph{over}-state yield, by $12\%$ here.

We note that the printed identity $\E[Y] = 1 + \sum_j \Prob[A_j]$ is a
bookkeeping check only: both sides are computed from the same run, so it verifies
our accounting rather than testing Proposition~\ref{prop:yield}, which is in any
case a definitional identity.

\subsection{The HMM identity and the impossibility of the bridge}

Theorem~\ref{thm:hmm} is verified against brute-force enumeration over all $M^T$
hidden paths, per $(t,j)$ rather than on aggregates of the final state, for $20$
random HMMs. A negative control confirms the hypothesis split: with a signed
transition matrix the sum-product identity still holds exactly (it is pure
algebra) while the max-product reading breaks, since $(\max,\times)$ distributes
only over $\R_{\ge0}$.

Theorem~\ref{thm:noembed} is confirmed directly. Under a gated linear-attention
update the set of rows on which two trajectories differ stays at $\{1\}$; under
the forward algorithm a single differing coordinate spreads to all of $[m]$ in
one step. The commutator of two observation-conditioned transition operators is
$\|\Psi_0\Psi_1 - \Psi_1\Psi_0\|_{\max} = 4.5\times10^{-2}$, so they are not
simultaneously diagonalisable. Both obstructions hold numerically.

\subsection{Capacity: the bound is never violated, achievability is hard to show}

Theorem~\ref{thm:recall} predicts that a $b$-bit interface can retain $n$ binary
values iff $b \ge n$. Sweeping both, no cell with $b < n$ ever reached the
success threshold --- the bound is never violated, which is the only direction in
which a lower bound can fail. The $n=2$ column reaches $1.00$ at $b=2,3,4$ as
predicted.

The achievability side is limited by optimisation rather than by capacity, and
this nearly misled us. An initial sweep at $800$ training steps produced seven
cells that failed \emph{with sufficient capacity}, including $b=4$, $n=2$, where
$16$ symbols must store $2$ bits. No capacity argument can explain that, since a
$4$-bit interface strictly contains a $3$-bit one and $b=3$, $n=2$ succeeded in
the same sweep; at $3000$ steps that cell reaches $1.000$ on every seed. An
undertrained discrete bottleneck is empirically indistinguishable from a capacity
ceiling, and a symmetric ``fraction of cells matching'' score obscures this by
weighting both error directions equally when only one of them can refute the
theorem. We report the violation count separately for this reason.
